\documentclass[10pt,a4paper]{amsart}
\usepackage{amsmath,amssymb,amsthm,graphicx}
\usepackage{verbatim}
\newenvironment{reference}{\comment}{\endcomment}
%\newenvironment{reference}{}{}
\newenvironment{slogan}{\comment}{\endcomment}
\newenvironment{history}{\comment}{\endcomment}

% For commutative diagrams we use Xy-pic
\usepackage[all]{xy}

% We use 2cell for 2-commutative diagrams.
\xyoption{2cell}
% Theorem environments.
%
\theoremstyle{plain}
\newtheorem{theorem}[subsection]{Theorem}
\newtheorem{proposition}[subsection]{Proposition}
\newtheorem{lemma}[subsection]{Lemma}

\theoremstyle{definition}
\newtheorem{definition}[subsection]{Definition}
\newtheorem{example}[subsection]{Example}
\newtheorem{exercise}[subsection]{Exercise}
\newtheorem{situation}[subsection]{Situation}

\theoremstyle{remark}
\newtheorem{remark}[subsection]{Remark}
\newtheorem{remarks}[subsection]{Remarks}

\numberwithin{equation}{subsection}

% Macros
%
\def\lim{\mathop{\mathrm{lim}}\nolimits}
\def\colim{\mathop{\mathrm{colim}}\nolimits}
\def\Spec{\mathop{\mathrm{Spec}}}
\def\Hom{\mathop{\mathrm{Hom}}\nolimits}
\def\Ext{\mathop{\mathrm{Ext}}\nolimits}
\def\SheafHom{\mathop{\mathcal{H}\!\mathit{om}}\nolimits}
\def\SheafExt{\mathop{\mathcal{E}\!\mathit{xt}}\nolimits}
\def\Sch{\mathit{Sch}}
\def\Mor{\mathop{\mathrm{Mor}}\nolimits}
\def\Ob{\mathop{\mathrm{Ob}}\nolimits}
\def\Sh{\mathop{\mathit{Sh}}\nolimits}
\def\NL{\mathop{N\!L}\nolimits}
\def\CH{\mathop{\mathrm{CH}}\nolimits}
\def\proetale{{pro\text{-}\acute{e}tale}}
\def\etale{{\acute{e}tale}}
\def\QCoh{\mathit{QCoh}}
\def\Ker{\mathop{\mathrm{Ker}}}
\def\Im{\mathop{\mathrm{Im}}}
\def\Coker{\mathop{\mathrm{Coker}}}
\def\Coim{\mathop{\mathrm{Coim}}}

% Boxtimes
%
\DeclareMathSymbol{\boxtimes}{\mathbin}{AMSa}{"02}

%
% Macros for moduli stacks/spaces
%
\def\QCohstack{\mathcal{QC}\!\mathit{oh}}
\def\Cohstack{\mathcal{C}\!\mathit{oh}}
\def\Spacesstack{\mathcal{S}\!\mathit{paces}}
\def\Quotfunctor{\mathrm{Quot}}
\def\Hilbfunctor{\mathrm{Hilb}}
\def\Curvesstack{\mathcal{C}\!\mathit{urves}}
\def\Polarizedstack{\mathcal{P}\!\mathit{olarized}}
\def\Complexesstack{\mathcal{C}\!\mathit{omplexes}}
% \Pic is the operator that assigns to X its picard group, usage \Pic(X)
% \Picardstack_{X/B} denotes the Picard stack of X over B
% \Picardfunctor_{X/B} denotes the Picard functor of X over B
\def\Pic{\mathop{\mathrm{Pic}}\nolimits}
\def\Picardstack{\mathcal{P}\!\mathit{ic}}
\def\Picardfunctor{\mathrm{Pic}}
\def\Deformationcategory{\mathcal{D}\!\mathit{ef}}
\usepackage{cleveref}
\crefname{originalreference}{original reference}{original references}
\Crefname{originalreference}{Original reference}{Original references}
\makeatletter
\newlabel{lemma-algebra-relations}{{0FHT}{source}{}{sourcepointer-1}{}}
\newlabel{lemma-algebra-relations@cref}{{[originalreference][0][]0FHT}{[1][0][]source}{}{}{}}
\makeatother
\crefname{originalreference}{original reference}{original references}
\Crefname{originalreference}{Original reference}{Original references}
\makeatletter
\newlabel{lemma-dim-0}{{0FHQ}{source}{}{sourcepointer-2}{}}
\newlabel{lemma-dim-0@cref}{{[originalreference][0][]0FHQ}{[1][0][]source}{}{}{}}
\makeatother
\crefname{originalreference}{original reference}{original references}
\Crefname{originalreference}{Original reference}{Original references}
\makeatletter
\newlabel{lemma-generated-by-separable}{{0FH5}{source}{}{sourcepointer-3}{}}
\newlabel{lemma-generated-by-separable@cref}{{[originalreference][0][]0FH5}{[1][0][]source}{}{}{}}
\makeatother
\title{Distinct cycle classes of rational points are not separated by divisibility}
\author{The Stacks Project Authors}
\date{stacks-project:weil.tex:local-snapshot-20261012; GNU FDL 1.2 or later}
\begin{document}
\maketitle
\noindent\textit{Editorial note.} All delivered mathematical text is the author's own, copied from the pinned Stacks Project chapter file weil.tex; no statement or proof step was rewritten, reordered or removed.\par\medskip
\noindent\textit{Editorial note.} The counted proof is the author's complete argument: the reduction to an irreducible scheme, the dualization of the cohomological relation into an identity of traces through the pairing axiom, and the application of the algebraic trace relation.\par\medskip
\noindent\textit{Editorial note.} The label lemma-relations-classes-points names the counted statement, and the proof environment immediately following it in the source is delivered verbatim as its proof.\par\medskip
\noindent\textit{Editorial note.} The counted statement is carried by the author's own lemma environment in the chapter source.\par\medskip
\noindent\textit{Editorial note.} Pre-open gate: python operations/agent-authoring/tools/precard\_gate.py --source weil.tex --label lemma-relations-classes-points --start 3655 --end 3694 --local-id 28141 returned RESULT=OK (L28141 free, occupancy FREE, public interval clear). Pre-assembly tag self-check: canonical tag resolved and not present in reference/public-items.jsonl.\par\medskip
\noindent\textit{Editorial note.} The statement and proof are copied from the pinned chapter file; the cited inputs are kept as external pointers with their canonical Stacks tags.\par\medskip
\noindent\textit{Allowed standard prerequisites (editorial).}\begin{itemize}
\item Generation of zero cycles by points with separable residue field. Same chapter, the lemma that zero cycles are generated by classes of closed points with separable residue fields. Allows writing a multiple of the difference of the two points as a combination of classes of points with separable residue fields.
\item Poincare duality for zero dimensional schemes and the trace description of the degree map. Same chapter, the computation of the cohomology of a zero dimensional scheme, including the identification of the degree map with the trace, and axiom (A)(b) expressing the pairing as pullbacks followed by traces. Dualizes the cohomological relation into a relation between trace maps out of the degree zero cohomology.
\item A relation between traces of surjective algebra maps forces equality. Same chapter, the algebraic lemma stating that a relation between traces of surjective maps of finite separable algebras forces the two surjections to agree. Concludes that the two pullback maps along the points coincide, contradicting the assumption on the cycle classes.
\end{itemize}\medskip
\ifdefined\hypertarget\hypertarget{sourcepointer-1}{}\fi
\noindent\textit{Source reference.} Original reference 0FHT: sources/stacks-weil/pinned-original-source.tex, source lines 3591--3591. Prerequisite pointer invoked at proof line(s) 3692; it names a standard fact of the same chapter that the counted argument uses, and every step of the counted proof built on it is retained verbatim in the delivered body.\par\medskip
\ifdefined\hypertarget\hypertarget{sourcepointer-2}{}\fi
\noindent\textit{Source reference.} Original reference 0FHQ: sources/stacks-weil/pinned-original-source.tex, source lines 3432--3432. Prerequisite pointer invoked at proof line(s) 3681; it names a standard fact of the same chapter that the counted argument uses, and every step of the counted proof built on it is retained verbatim in the delivered body.\par\medskip
\ifdefined\hypertarget\hypertarget{sourcepointer-3}{}\fi
\noindent\textit{Source reference.} Original reference 0FH5: sources/stacks-weil/pinned-original-source.tex, source lines 2233--2233. Prerequisite pointer invoked at proof line(s) 3672; it names a standard fact of the same chapter that the counted argument uses, and every step of the counted proof built on it is retained verbatim in the delivered body.\par\medskip
\begin{lemma}
\label{lemma-relations-classes-points}
Assume given (D0), (D1), (D2), and (D3) satisfying (A), (B), and (C).
Let $k'/k$ be a finite separable extension.
Let $X$ be a smooth projective scheme over $k'$.
Let $x, x' \in X$ be $k'$-rational points.
If $\gamma(x) \not = \gamma(x')$, then
$[x] - [x']$ is not divisible by any integer $n > 1$ in $\CH_0(X)$.
\end{lemma}

\begin{proof}
If $x$ and $x'$ lie on distinct irreducible components of $X$, then
the result is obvious. Thus we may $X$ irreducible of dimension $d$.
Say $[x] - [x']$ is divisible by $n > 1$ in $\CH_0(X)$.
We may write $[x] - [x'] = n(\sum m_i [x_i])$ in $\CH_0(X)$
for some $x_i \in X$ closed points
whose residue fields are separable over $k$ by
Lemma \ref{lemma-generated-by-separable}.
Then
$$
\gamma([x]) - \gamma([x']) = n (\sum m_i \gamma([x_i]))
$$
in $H^{2d}(X)(d)$. Denote $i^*, (i')^*, i_i^*$ the pullback maps
$H^0(X) \to H^0(x)$, $H^0(X) \to H^0(x')$, $H^0(X) \to H^0(x_i)$.
Recall that $H^0(x)$ is a finite separable $F$-algebra
and that $\int_x : H^0(x) \to F$ is the trace map
(Lemma \ref{lemma-dim-0}) which we will denote $\text{Tr}_x$.
Similarly for $x'$ and $x_i$. Then by Poincar\'e duality in the form of
axiom (A)(b) the equation above is dual to
$$
\text{Tr}_x \circ i^* - \text{Tr}_{x'} \circ (i')^* =
n(\sum m_i \text{Tr}_{x_i} \circ i_i^*)
$$
which takes place in $\Hom_F(H^0(X), F)$. Finally, observe that
$i^*$ and $(i')^*$ are surjective as $x$ and $x'$ are $k'$-rational
points and hence the compositions $H^0(\Spec(k')) \to H^0(X) \to H^0(x)$
and $H^0(\Spec(k')) \to H^0(X) \to H^0(x')$ are isomorphisms.
By Lemma \ref{lemma-algebra-relations} we conclude that $i^* = (i')^*$
which contradicts the assumption that $\gamma([x]) \not = \gamma([x'])$.
\end{proof}

\end{document}
